% Generated draft: requires mathematical and source review.
\documentclass[11pt]{article}
\usepackage[margin=25mm]{geometry}
\usepackage{fontspec}
\setmainfont{Latin Modern Roman}
\usepackage{amsmath,amssymb,mathtools,unicode-math}
\setmathfont{Latin Modern Math}
\usepackage{enumitem,needspace,xurl,hyperref}
\setlength{\parindent}{0pt}
\setlength{\parskip}{5pt}
\setlength{\emergencystretch}{3em}
\newcommand{\N}{\mathbb{N}}
\newcommand{\Z}{\mathbb{Z}}
\newcommand{\Q}{\mathbb{Q}}
\newcommand{\R}{\mathbb{R}}
\newcommand{\C}{\mathbb{C}}
\newcommand{\F}{\mathbb{F}}
\newcommand{\A}{\mathbb{A}}
\newcommand{\PP}{\mathbb{P}}
\newcommand{\E}{\mathbb{E}}
\newcommand{\eps}{\varepsilon}
\newcommand{\dd}{\,\mathrm{d}}
\begin{document}
\section*{UnsolvedMath ID: 20001892}
\textbf{Source problem code:} catalogue entry
\subsection{Definitions and mathematical statement}
The source's Question \(5.4.1\) asks whether there exists a closed negatively curved Riemannian manifold \((M,g)\), of real dimension greater than two, whose geodesic flow has \(C^1\) stable and unstable bundles \(E^s,E^u\), while \(g\) is not \(1/4\)-pinched. Here the Anosov splitting on \(SM\) is \(T(SM)=E^s\oplus E^0\oplus E^u\), with \(E^0\) tangent to the flow; \(C^1\) refers to these bundles as distributions on all of \(SM\).

The pinching endpoint is essential. For the global convention, put
\[
 k_{\min}=\min_{x,\sigma}(-K_g(\sigma))>0,\qquad
 k_{\max}=\max_{x,\sigma}(-K_g(\sigma)).
\]
Weak quarter pinching means \(k_{\max}/k_{\min}\leq4\), whereas strict quarter pinching means \(k_{\max}/k_{\min}<4\). Under the weak interpretation, the requested metric must have ratio greater than \(4\); the aid reports that no such example with both bundles \(C^1\) was found in literature checked through 8 August 2026. This is a dated open-status qualification, not a nonexistence theorem.

Under the strict interpretation, compact complex-hyperbolic quotients of complex dimension at least two already give the reported positive answer: their curvature ranges from \(-4\) to \(-1\), the ratio is exactly \(4\), and the locally symmetric stable and unstable bundles are smooth. They fail strict quarter pinching but satisfy weak quarter pinching, so they do not answer the intended beyond-weak-endpoint question.

The appended source discussion motivates a perturbation program from [Has-Wil1999], Section \(4\), especially Proposition \(12\), seeking open metric sets with low-regularity Anosov splitting on large subsets. It requires a suitable base with widely separated unstable expansion rates and says that suitable geodesic-flow bases were not then known. It also points to Pollicott's question in the next subsection, whose statement is not part of this record.

For the aid's exact parallel two-rate symplectic Anosov model, let the slow and fast positive rates be \(0<a\leq B\), with paired negative stable rates. ``Exact parallel'' means the invariant rate subbundles have the prescribed exponential rates without variable-rate errors. The aid's reduction of the Hasselblatt--Wilkinson \(\alpha\)-u-spread parameter feasibility is
\[
 \text{the specified spread inequalities are feasible}
       \quad\Longleftrightarrow\quad \alpha B>2a.
\]
The full definition and auxiliary parameter inequalities of \(\alpha\)-u-spread are not supplied, so this equivalence is retained as a reported reduction in that exact model. Complex hyperbolic space has parallel curvature sectors \(-1,-4\) and positive rates \(a=1,B=2\). Thus it fails the strict \(1\)-u-spread condition at equality, but satisfies the \((1+\eps)\)-version for every \(\eps>0\). This endpoint obstruction does not prove that no other perturbation argument works or that a requested beyond-quarter-pinched example cannot exist.
\subsection{Short English statement}
The question depends on whether quarter pinching includes equality. Complex-hyperbolic examples answer the strict version, while the beyond-weak-bound version is reported open through 8 August 2026. The aid's two-rate reduction explains why complex-hyperbolic rates \(1,2\) sit exactly at, and fail, the strict \(1\)-u-spread threshold.
\subsection{Sources}
\begin{itemize}
\item Local UnsolvedMath registry, permanent ID \(20001892\), source code \texttt{catalogue entry}; canonical collection \texttt{aim-geometry-notes.json}, notes[229]. Source-item field \(5.4\), actual question \(5.4.1\).
\item Keith Burns and Vladimir Matveev, ``Open problems and questions about geodesics,'' AIM workshop ``Geodesics'': \url{https://aimath.org/pastworkshops/geodesicsproblems.pdf}. The supplied verification reports inspection from the beginning of Section \(5.4\) through the beginning of \(5.4.2\), with only OCR repairs to the mathematical question. This drafting session did not open the PDF.
\item Hasselblatt--Wilkinson, supplied label [Has-Wil1999], Section \(4\), Proposition \(12\) and end-of-section discussion, for the low-regularity perturbation program. The source names Pollicott as further motivation and describes complex-hyperbolic exponents \(1,2\). Full reference entries and the following Pollicott question are not supplied.
\item Supplied research attempt \(1\), \texttt{partial\_result}, run status ``open,'' rigor label ``proved,'' novelty confidence ``low.'' Registry status \texttt{partially\_solved}, classification \texttt{PARTIAL-PROGRESS}, updated 14 August 2026. Its literature qualification is explicitly checked only through 8 August 2026. Original remarks and literature fields are empty; no proposer or year is supplied.
\end{itemize}
\end{document}
